% drivers/project-notes.tex
\documentclass{tex-notes}

\title{\vspace{-2em}Proyecto Final: DeepLearning\\
  {\large \textnormal{%
    $G_\theta : \mathbb{R}^{128} \longrightarrow [-1, 1]^{1 \times 64 \times 64}$%
  }}}
\author{}
\date{}

\begin{document}
\maketitle
\vspace{-2.5em}

%% =====================================================================
%% Per-block walkthrough
%% =====================================================================

\input{content/proyecto/00-introduccion}
dominio de aplicación
\input{content/proyecto/00-problema-justificacion}
descripción del problema
justificación técnica

El problema está bien definido
La aplicación generativa es clara
En el dominio elegido tiene sentido y relevancia

\input{content/proyecto/00-datos-metodos}

datos utilizados
propuesta de arquitectura
La arquitectura elegida es adecuada
La propuesta está bien explicada
Existe una justificación técnica sólida basada en los temas del curso
estrategia de entrenamiento

\input{content/proyecto/00-implementacion}
limitaciones, riesgos y trabajo futuro
El sistema fue implementado correctamente
El proyecto funciona de manera coherente
El demo refleja la solución propuesta

El demo funciona correctamente
Es claro y usable
Permite apreciar el valor del proyecto

El código está ordenado
El repositorio es entendible
El README y las instrucciones permiten reproducir o comprender la solución

%% =====================================================================
%% Closing sections
%% =====================================================================

\input{content/proyecto/00-resultados-discusion}
evaluación de resultados
limitaciones, riesgos y trabajo futuro
\input{content/proyecto/00-referencias}
referencias

\input{content/proyecto/00-script-video}
El video comunica bien la solución
Presenta claramente la arquitectura, demo y resultados
Se mantiene dentro del tiempo solicitado

\end{document}
